\usepackage[spanish]{babel}
%!TEX root = ../main.tex
%!LTeX: language=es, en

\chapter{Un rápido diagnóstico}

\section{Conocimiento e impresiones adecuadas}

%Si Una vez que determinamos el valor de verdad de una oración---el proceso en el
%cual juzgamos si una afirmación es verdadera o falsa--- criterio de equivalencia de
%juicios es el cambio de nombres 

%Los condicionales contrafácticos son parte de la caja de herramientas de cualquier
%investigador contemporáneo. Sin embargo, 

\noindent No debe haber mucha resistencia si afirmo que la investigación científica
tiene como objetivo alguna forma de conocimiento: ya sea \textit{predecir} fenómenos,
\textit{explicar} fenómenos, \textit{catalogar} fenómenos, etc. En principio, parece
sensato también decir que al realizar esta actividad y lograr sus objetivos, los
investigadores descubren cómo \textit{de hecho} es el mundo. Otras actividades
humanas (la música, la cocina o la pintura) no tienen claramente objetivos
epistémicos.

Asumiré que la fricción no aumenta si afirmo que si dichos objetivos se satisfacen,
de alguna forma obtenemos \textit{conocimiento.} Dicho de esta forma, se sigue que si
los investigadores logran cualquiera de sus objetivos, su actividad de alguna manera
debería resultar en una \textit{creencia,} \textit{verdadera} y
\textit{justificada.}\footnote{ Recordemos que los contraejemplos Gettier muestran
que las condiciones anteriores, en conjunto no son suficientes
\parencite{gettier1963} para tener conocimiento, esto no implica que los elementos
involucrados no sean necesarios.}

Siendo más explícito, es necesario que: si un sujeto cualquiera tiene conocimiento de
que $p$, entonces $p$ es una \textit{proposición} verdadera. Si esto es correcto para
todo objetivo epistémico, entonces vale para los casos particulares en los que el
sujeto en cuestión tenga como objetivo realizar una \textit{predicción,} o una
\textit{explicación} o una \textit{catalogación} de fenómenos. En términos crasos,
los investigadores obtienen como resultado una \textit{impresión} verdadera de cómo
es el mundo.

Para enfatizar por qué es relevante obtener una impresión adecuada de cualquier fenómeno en
cuestión, basta con notar que cuando los investigadores fallan en sus objetivos, el
impacto se extiende hasta dominios no epistémicos. El caso más claro de esto es que
usamos los resultados de la investigación, por ejemplo, para informar políticas
públicas para lidiar con el cambio climático y el calentamiento global
\parencite{stern2007}. Ambos son fenómenos que involucran aspectos económicos y
sociales, tanto como químicos, físicos y biológicos. Este fenómeno es tan apremiante
que dar acertadamente (en términos de conocimiento) con la respuesta
\textit{correcta}, es de suma importancia. Si bien no siempre logramos el objetivo en
cuestión, al menos podemos acordar en que las políticas públicas diseñadas para
resolver problemas sociales \textit{deberían} estar diseñadas tomando en cuenta la
evidencia pertinente según sea el caso.

Esto, en espíritu, es algo que deja bien ilustrado Lynch al decir ``Esta es la razón
por la cuál las democracias alrededor del mundo deben redoblar su compromiso con la
infraestructura epistémica---aquellas instituciones y prácticas que está a la base de
la búsqueda de la verdad. Esto significa proteger a la ciencia y a la historia de
intrusión política''\footnote{ \say{That is why democracies around the world must
urgently redouble their commitment to epistemic infrastructure---to the institutions
and practices that undergird the responsible search for truth. That means protecting
science and history from political intrusion}} \parencite{lynch2025truth}.

Estoy sólo parcialmente de acuerdo con lo que dice Lynch---resta discutir a qué se
refiere con \say{democracia} e \say{intrusión política}---hasta el punto en el no es
controvertido decir que al menos una parte de nuestra política pública debería basar
sus decisiones en la infraestructura epistémica responsable de la búsqueda de la
verdad. Esta infraestructura, según entiendo, es la actividad que realizan aquellos
que se dedican a la investigación científica.

De suyo, este traslape es motivación suficiente para analizar el papel que la
\textit{verdad} juega en la investigación, pero debo aclarar que esto sólo empuja el
problema un paso atrás. Si bien puedo asegurar que vagamente acordamos con lo
mencionado en los párrafos anteriores, esta imagen intuitiva tiene problemas que los
filósofos de la ciencia han discutido ampliamente. Es en esta disciplina y afines
donde encontramos razones para dudar que la investigación científica---la
infraestructura epistémica encargada de la búsqueda de la verdad---sea no sólo
exitosa en sus objetivos, sino si el objetivo es de hecho el adecuado.

Razones para dudar de que la verdad esté entre los objetivos de la investigación
científica, van desde argumentos que defienden que las creencias falsas son y
pueden ser de utilidad para lograr objetivos epistémicos \parencite{elgin2004,
elgin2017true}. Que la investigación científica no produce verdad en ningún sentido
interesante \parencite{Shapin1995}. Que aún si los objetivos epistémicos son valiosos
porque implican verdad, hay otros estados epistémicos---como la comprensión, la
sabiduría, la racionalidad, etc---que no implican verdad\footnote{Esto es que si la
ciencia tiene objetivos epistémicos, entonces cabe la posibilidad de que la ciencia
logre sus objetivos sin que haya ningún sentido de \textit{verdad} involucrado} y son
estados epistémicos valiosos \parencite{kvanvig2003}.

El desarrollo de los argumentos anteriores involucra detalles que dejaré para
capítulos posteriores, pero puedo asegurar que las afirmaciones previas concuerdan en
al menos un punto: que podemos deshacernos de la noción de \textit{verdad} en
investigación. Podemos deshacernos de ella porque no es útil para analizar el
desarrollo de la investigación científica.

\paragraph{Me parece que} afirmaciones como las de Elgin, Shapin y afines, resultan
de la mala concepción del concepto de ``verdad'' que suele estar involucrado en estos
debates. Debo señalar que una parte de mi trabajo---y una de las premisas para mi
argumento---es intentar convencerlos de que hay que recuperar un concepto particular
de ``verdad,'' que es el que me permite aclarar parte del debate. Entonces, para
comenzar a trabajar, viene bien definir de manera provisoria algunas de las
condiciones que por lo menos debería cumplir el concepto de ``verdad'' que quiero
recuperar. PAra lograr esto, a lo largo de este trabajo voy a hablar en términos de un
marco \textit{proposicional.} Las proposiciones son los objetos que pueden ser
verdaderos o falsos. El valor de verdad de una proposición está determinado por el
contenido semántico de la misma. Contenido que depende de una relación externa entre
las proposiciones y el mundo. Esa última condición es suficiente para decir que la
teoría de la verdad en la cuál estoy pensando, cae bajo la categoría de una ``teoría
correspondentista'' de la verdad \cite{rasmussen2014, sep-truth}. Es
correspondentista porque la tesis central es que sea cuál sea el contenido de una
oración---y dada una proposición cualquiera---su valor de verdad está determinado por
la relación que hay entre la proposición y el mundo, esto es, que la oración
corresponde adecuadamente con ``los hechos.''

\section{Límites y alcances}

\noindent Quiero recordar al lector que lo anterior es una caracterización, y que no debemos
confundir esto con una teoría sustantiva de la verdad \parencite{engel2014truth}.
Precisar y defender una teoría sustantiva, es un análisis que está fuera de los
propósitos de este trabajo. No me corresponde defender una teoría correspondentista
particular. Quiero enfatizar además que no he mencionado algo acerca de cuál es la naturaleza de las
proposiciones \parencite{sep-propositions}, ni acerca de la naturaleza de la relación
entre las proposiciones y el mundo \parencite{ramsey1927facts}. Tampoco he afirmado
algo en torno a cómo, independiente a la naturaleza de las proposiciones, es que las
proposiciones se relacionan con nuestro aparato cognitivo
\parencite{frege1918thought}. Y no he señalado nada en torno a la naturaleza de los
hechos a los cuáles las proposiciones corresponden \parencite{austin1950truth}.

Detallar cualquiera de estos aspectos implica diferentes versiones de una teoría de
la verdad. Además, siguiendo la taxonomía que ofrece Engel
\parencite{engel2014truth}, y términos estándar de la literatura filosófica, tampoco
me he comprometido con qué son los \textit{truth-bearers}, ni los
\textit{truth-makers} \parencite{sep-truth}. Sólo me he comprometido con que es una
relación en la que el mundo juega un papel.

Afortunadamente, creo que lo único que necesito para completar este trabajo es asumir
que \textit{el mundo juega un papel}, porque este trabajo no versa sobre la
justificación de una teoría de la verdad, sino de la evaluación de juicios en la
práctica científica, especialmente en los casos que afirmamos que los juicios son
verdaderos o falsos\footnote{Aquí estoy siguiendo a Frege, hasta el punto en el que
afirma que las ciencias empíricas no solamente conciernen al sentido, sino a la
referencia \cite{fre}. Lo que significa, en términos crasos, es que entender que una
oración tiene contenido \textit{no es suficiente} para juzgar si la oración es
verdadera o falsa. Una vez que detectamos que una oración tiene contenido, juzgarla
es pasar a reconocer---investigar es la palabra que se antoja---que hay un nombre del
cual se afirma que tiene una propiedad. Juzgar, en un contexto cotidiano, es afirmar
de la oración ``mi cafetera es negra'' que es verdadera o falsa. Juzgo a la oración
verdadera si busco en mi cocina y mi cafetera es negra. También estoy siguiendo al
autor alemán en que juzgar a la oración como falsa, implica una disyunción: la
negación de la oración anterior es o bien que ``mi cafetera no es negra,'' o bien que
mi cafetera no existe. En el caso anterior, es claro que las condiciones en la que
diríamos que \textit{miento} sobre el color de mi cafetera es cuando o bien nunca
tuve una cafetera---fue una ilusión, nunca fue mía a pesar de que pensé que sí lo
era, ya no tengo cafetera, etc,---o bien mi cafetera es de un color distinto. Estoy
deliberadamente usando \textit{mentir}, porque, si bien este es un ejemplo de un
juicio en la vida cotidiana, no veo cómo difiere del contexto de la investigación
científica, por lo menos en lo que respecta al proceso de juzgar. Si esto es verdad y el contexto de investigación cumple con las suficientes características para caer bajo el principio de cooperación de Grice \parencite{}, entonces podemos decir que la investigación es una actividad en
la que esperamos que los involucrados, por lo menos, \textit{no mientan}---acorde con
la máxima de contenido y relevancia que discute el autor británico. Pero también es claro que es un contexto mas
rígido que el contexto de la vida cotidiana, porque esperamos que las personas
involucradas cumplan ciertos estándares al momento de juzgar si una afirmación es
verdadera o falsa. Un estándar obvio es que al menos haya intentado juzgar a sus oraciones de manera tal que pueda afirmar si su oración es verdadera o falsa. Que afirma de un objeto que es de tal o cual manera. Asumimos que la disyunción discutida por Frege no se presenta, porque asumimos que cuando una persona afirma que ``los electrones tienen carga negativa,'' asumimos que ambos entendemos, aunque sea vagamente, qué es un electrón y qué es tener carga negativa. La pregunta de si el electrón existe o no, no se presenta, porque asumimos que estamos hablando de algo, poniendolo en forma de pregunta ¿por que a alguien en el contexto de investigación le interesaría hablar de objetos que no existen? Es excluyente con la actividad de juzgar caracterizada anteriormente porque, pensando en objetos que caen bajo el dominio de las entidades ficticias, como afirma Frege, perdemos el punto del goce estético cuando hacemos juicios sobre si lo que sucede en las novelas de Conan Doyle es verdadero o falso. En segundo lugar  porque También estoy
asumiendo, de forma muy vaga, que en el contexto de la investigación es donde se
presenta la pregunta de si los juicios son verdaderos o falsos y del hecho de que la
investigación sea una actividad colectiva, se sigue que en el contexto de
investigación, en donde se presenta la pregunta acerca de los juicios esto es un
proposiciones, y además estoy sugiriendo que si la investigación, sea de cualquier
tipo, le conciernen los juicios de oraciones verdaderas o falsas }. En los casos más
obvios, juicios de la vida cotidiana sobre la evaluación de \textit{juicios
contrafácticos.} Lo que quiero mostrar es que la verdad \textit{sí importa} cuando
\textit{evaluamos} si un \textit{condicional contrafáctico} es verdadero o falso. Es
la parte de la evaluación de estos condicionales de lo que trata este trabajo ¿bajo
qué criterios decimos que un condicional contrafáctico es verdadero? 

En principio, es extraño decir que el mundo juega un papel dado que los condiconales
contrafácticos son contrarios a los hechos. Estos condicionales suelen ser
subjuntivos y es estándar asumir que cuando afirmamos un condicional de este estilo,
por ejemplo ``si hubiera salido 10 minutos antes, entonces hubiera llegado puntual al
trabajo,'' lo juzgamos bajo el supuesto de que el antecedente es falso. Porque está afirmación es natural cuando de hecho no salí 10 minutos antes y llegué tarde al trabajo. Si los condicionales contrafácticos compartieran la semántica del condicional material, por defecto todos los contrafácticos serían verdaderos. Por utilizar
téminos de Frege, 

en el mercado filosófico, tengo a la mano una formulación particular
que me permite analizar los aspectos de la investigación científica que me importan
en este trabajo que son \textit{verdad} y \textit{justificación.} De esto es de lo
que trata el primer capítulo.

\paragraph{Estoy suponiendo} además que la investigación científica es una
\textit{actividad} humana--de hecho es un trabajo por el cuál se paga.
Siendo este el caso, es más que obvio que el trabajo de los investigadores
está atravesado por decisiones políticas y económicas. Los sesgos y
las creencias particulares de las personas dedicadas a este trabajo
influyen en cómo desarrollan su actividad. La investigación también es una
actividad colaborativa. Nadie sostendría que hacer investigación (en
cualquier ámbito) es una actividad que puedo realizar haciendo un
soliloquio en mi casa. Para quienes digan lo contrario, la respuesta obvia
es que no debería hacer una tesis para obtener un grado: debería lograrlo
sin tener que presentar un escrito.\footnote{Para aquellos que
respondadn que esto es sólo por razones institucionales, sólo puedo decir
que es una institución que al menos es tan antigua como Arquímedes---quien
no publicó el trabajo que realizó en ``El método,'' porque no le quedaba
claro que los resultados obtenidos fueran lo suficientementes sólidos para
ofrecer una justificación adecuada para sus pares
%\cite{archimedes_method_1912}.
} El hecho de que parte de la investigación
consista en publicar resultaddos, es evidencia suficiente para concluir que
la colaboración entre invetigadores es una de las características
esenciales de este trabajo.

La caracterización de la verdad que di más arriba deben bastar por ahora, ya que lo que mi trabajo aporta a la literatura depende de dos premisas que sí voy a defender, a saber, que \textit{el objetivo de la investigación científica es alguna forma de conocimiento} y que \textit{la verdad es condición necesaria para tener conocimiento}. Debo también mencionar que espero que el lector conceda la suposición de que la investigación es una actividad colaborativa. Si la teoría de la verdad es correspondentista y, si en principio las personas involucradas en este trabajo tienen el interés de hablar con verdad---o al menos evitar los errores en la medida de lo posible,---de no ser así, las consecuencias epistémicas y prácticas son repugnantes\footnote{Creo que esto incluye el caso en el que las hipótesis que sostuvieron los investigadores fueron---o son,---falsas o inadecuadas.}

Los que suele generar incomodidades entre filósofos de la ciencia es, me parece, que de alguna manera las afirmaciones anteriores no pueden ser todas verdaderas al mismo tiempo.

\begin{enumerate}
\item El resultado de la investigación científica es el conocimiento
\item La investigación científica es objetiva.
\item La investigación científica es colaborativa.
\item El conocimienbto es valioso porque es verdadero.
\item La investigación científica tiene influencias y consecuencias prácticas (políticas, económicas, etc.)
\item La verdad es una relación externa independiente de las creencias humanas.
\end{enumerate}

Parece un supuesto generalizado que si la investigación científica necesariamente está involucrada en la política pública, entonces no puede ser objetiva. Si no es objetiva, \textit{a fortiori} no es verdadera.  Aún si quisiéramos recuperar la noción de objetividad, argumentando que el concepto de \textit{objetividad} no está %emparentado con el de \verdad{verdad} \parencite{truthwoObjectivity}, los sesgos políticos económicos (y de cualquier otro tipo), minarían incluso este intento.

Es más que obvio que lo anterior es preocupante, más si insisto en que Lo anterior es preocupante, especialmente tomando en consideración los argumentos de Elgin, Kvanvig o Shapin. La cuestion semántica





\section{¿Qué podemos aprender de la historia de la ciencia?}
%\label{sub:aprender}

\noindent Quiero comenzar esta sección confesando lo siguiente: estoy de
acuerdo con la autora, en particular con el carácter de la conclusión
anterior, esto es, que es complicado--si no es que imposible--, separar
entre los procesos externos y los procesos internos que influyen en la
investigación científica. Pero aceptar esto --eso quiero argumentar-- no
implica que debamos deshacernos de la distinción entre el \emph{contexto de
investigación} y el \emph{contexto de descubrimiento}. En lo que resta de
la sección quiero hacer un breve repaso de los argumentos de la autora.

Yturbe nos dice que el contexto de descubrimiento, se dedica a analizar
factores externos que fueron influyendo en los factores internos; mientras
que el \emph{contexto de justificación} sólo se dedica a analizar los
factores internos al desarrollo teórico. Lo que ella llama \say{the
doctrine of the two contexts}\footnote{ \say{one of the philosophical
theses in favor of the contraposition between the internalist approach and
the externalist approach is the so-called doctrine of the two contexts,
%developed by positivism.} \cite[p. 75]{Yturbe1995} } La autora nos dice que
la doctrina de los dos contextos hace una distinción que no se puede hacer.
El contexto de justificación es siempre influenciado por el contexto de
descubrimiento porque la ideología, la economía, las relaciones de poder,
etc. influyen en la toma de decisiones\footnote{Por usar una caricatura: a
quién se le asigna presupuesto}, por lo tanto la doctrina es incorrecta.

Como dije al principio de esta sección, estoy de acuerdo con que es casi
imposible hacer la distinción entre \emph{historia interna} e
\emph{historia externa}, pero aceptar esto, no implica deshacernos de la
distinción \emph{contexto de descubrimiento} y \emph{contexto de
justificación}. Para sustentar esto, dependo de la caracterización que hace
la autora sobre \say{la doctrina de los dos contextos}, la caracterización
de la doctrina está ligeramente sesgada, si no es que completamente
inadecuada. En la siguiente sección quiero ofrecer mis razones para esta
conclusión.

\subsection{La doctrina}

%\noindent En el texto de \textcite[p.75]{Yturbe1995}, la autora ofrece una
descripción de la doctrina, ella señala que los filósofos doctrinarios
afirman que \say{according to this position, [factores externos], not only
fails to increase our understanding of scientific development, but even
obscures the fundamental question of the rational validation of scientific
arguments.} Además que \say{This conception constitutes the dominant
framework in which the philosophy of science has been developed, and
consists in drawing a radical distinction between the context of discovery
and the context of validation.}


Su argumento para decir que hay un colapso entre factores externos
(contexto de descubrimiento) y factores internos (contexto de
justificación) se encuentra en la página 85

\begin{quote}

	External factors are not only found in the context of discovery, they
	are present also in the development of the concepts, problems, methods,
	problem fields, etc. of scientific discourses: that is, external
	factors are present in the context of validation itself. There are
	scientific discourses in which ideological conceptions pass on to form
	part of the body of the science itself, functioning as principles which
	define its field of study or guide its research; thus, external factors
	can become internal.

\end{quote}

La doctrina de los dos contextos implica que podemos claramente separar
entre el \emph{contexto de descubrimiento} y el \emph{contexto de
justificación}. Pero los factores externos están presentes en el contexto
de justificación, tanto así que pueden convertirse en factores internos.
Esto implica que los contextos no son claramente separables, y que, por
tanto, la doctrina es incorrecta.

Contrario a lo que dice la autora

\begin{quote} The search for new explanatory programs is characterized
above all by the attempt to reconcile the internalist and externalist
approaches. But, in view of the fact that both approaches are committed to
theses concerning the nature of science that are not only incompatible, but
are unsustainable, their union without important changes in their
philosophical presuppositions cannot result in a viable program.
\parencite[p. 79]{Yturbe1995} \end{quote}

Creo que es verdad que no podemos distinguir entre factores externos y
factores internos en la ciencia. Sin embargo, me parece que lo que expresa
la cita anterior es claramente erróneo: no se siguie la conclusión \say{que
no pueden resultar en un programa viable}, porque depende de señalar que la
distinción entre \emph{contexto de descubrimiento} $|$ \emph{contexto de
justificación} es equivalente a la distinción \emph{factores externos} $|$
\emph{factores internos} y estas distinciones no son equivalentes.



\subsection{En defensa de la doctrina}

\noindent En años recientes ha habido un creciente interés en discutir las
%publicaciones del Círculo de Vienna \parencite{Bentley2023, Richardson2023,
%Suarez2024, Riel2014}. Si hago una tosca generalización, diría que
literatura reciente se ha dedicado a señalar que los miembros del Círculo
de Vienna, no eran tan herméticos como se pensaba. Además, la cantidad de
temas que discutieron es más vasta de lo que se creía. Particularmente
pensando en que cada uno de los miembros difería de los otros en tesis
centrales.

%\textcite{Suarez2024}, por ejemplo, discute que una preocupación genuina
sobre los modelos en la ciencia, precede, data y prosigue a los años del
Círculo de Vienna, el autor nos dice \say{I focus particularly on the
nineteenth-century modelers and summarize their insights and contributions,
which, I claim, remain essentially unsurpassed.} (p. 20)

Además, los miembros del Círculo de Vienna tuvieron una preocupación
genuina por las prácticas científicas--Neurath, por ejemplo-- Tenían
presente que ciertos factores externos pueden influenciar a la
%investigación científica. \textcite[p. 24]{Bentley2023} lo expresa mejor:
\say{Frank's historical work, which frequently anticipates Kuhnian or
post-Kuhnian themes, consistently emphasizes the significance of
non-scientific, external factors on the decision making of scientists.} Los
miembros del Círculo de Vienna tenían claro que los factores externos
influyen en el proceso de la investigación científica, presentes incluso en
el \emph{contexto de justificación}.

Hablando en particular del trabajo de Neurath y las tesis que sostenía,
Otto defendía una teoría coherentista de la justificación. Neurath
argumentaba que nuestras teorías científicas están subdeterminadas
empíricamente. En cualquier momento del tiempo hay hipótesis en pugna. Pero
si las teorías están empíricamente subdeterminadas, es posible que dos
teorías internamente coherentes y externamente inconsistentes entre sí,
convivan al mismo tiempo.

Pero si el escenario anterior es plausible, entonces es patente que aquél
que defienda una teoría coherentista de la confirmación, no puede resolver
el problema de la elección racional de teorías. Debido a que todo el tiempo
hay teorías en pugna, como filósofos, deberíamos poder explicar por qué es
racional elegir entre una teoría y sus rivales. Para poder resolver este
problema, Neurath defendía que los factores externos son información clave
para explicar la racionalidad de la elección de los investigadores.
Factores externos como los eventos históricos fortuitos.

%\textcite{Bentley2023} nos recuerda que Neurath fue falibilista sobre el
conocimiento, es decir, sostenía que cada una de nuestras creencias puede
estar equivocada. Sumado a esto, sostenía que la confirmación es holista:
no son oraciones particulares, sino grupos de oraciones, los que
contrastamos con el mundo. Bentley lo expresa mejor \say{$\ldots$ it is
always possible for a system of beliefs to be altered to accommodate a
statement or for the statement to be rejected $\ldots$} (p. 21).


%Esto último inició como un abstract que mandé a un taller y se fue
volviendo más largo mientras leía. %Hablando de modelos en la filosofía de
la ciencia durante el periodo del Círculo de Vienna, una figura importante
que mencionar es Mary Hesse. %Mary Hesse trabajó con especial atención el
uso de modelos en la ciencia.


%This historical chapter introduces the emergence in the nineteenth century
of what I call the modeling attitude. This is a stance toward scientific
work and discovery, and it continues to this day. (p. 43)
